\section{Conclusion}
\label{sec:conclusion}

\begin{revpar}
We presented \methodname{}, which adapts the SAM--CLIP correlation paradigm to
binary SOD with a fine-tuned CLIP ViT-L/14 encoder, a fixed prompt bank, and an
adaptive-fallback pseudo-prompt generator. It obtains the best S-measure on four of
five RGB benchmarks among the compared methods and, with RGB input only,
outperforms the listed RGB-D methods on SIP. Our analysis shows that the prompt
residual restores a gradient path without improving accuracy, and that prompt
design matters far less than adapting the visual encoder.

\noindent\textbf{Limitations.} Several effects are within seed variation, and only
part of the study uses three seeds. Published F-measure variants differ across
sources, and the SIP comparison uses different training data than the RGB-D
methods.
\end{revpar}
%% OLD: "three key innovations: a prompt-residual connection ..., a learned per-pixel
%% attention mechanism ..., and an adaptive-fallback PPG ... (1) the aggregation
%% mechanism matters more than prompt wording ... (2) architectural simplicity with
%% smart prompt handling outperforms complex decoders with fixed aggregation ...
%% achieves state-of-the-art S-measure on five RGB benchmarks and surpasses
%% depth-based methods on SIP using RGB input only."
